% ---
% title: Two exponential random variables
% tags: joint-continuous, marginal, exponential
% ---

\begin{problem}
	The random variables $X$ and $Y$ have joint density
	\[
		f(x,y)=\begin{cases}
			2e^{-x}e^{-2y} & x>0,\ y>0, \\
			0              & \text{otherwise.}
		\end{cases}
	\]
	(a) Find $P(X>1,\,Y<1)$.
	(b) Find the marginal densities $f_X$ and $f_Y$.
	(c) Find $P(X<Y)$.
\end{problem}

\begin{solution}
	\idea The probability of an event $\{(X,Y)\in A\}$ is the double integral
	of the joint density over the region $A$. Compute the inner integral first,
	holding the other variable fixed. A marginal density comes from integrating
	the joint density over the \emph{other} variable.

	\step{Check.} The density is never negative. Its total integral is $1$:
	\[
		\int_0^\infty\!\!\int_0^\infty 2e^{-x}e^{-2y}\,dx\,dy
		=\int_0^\infty e^{-x}\,dx\cdot\int_0^\infty 2e^{-2y}\,dy
		=1\cdot 1=1.
	\]
	(Each factor is the total integral of an exponential density.)

	\step{(a)} The event is $x>1$ and $0<y<1$ (the density is $0$ for
	$y\le0$, so the lower bound of $y$ is $0$). This region is a rectangle
	(an infinite strip to the right of $x=1$), so the bounds of $x$ do not
	depend on $y$:
	\[
		P(X>1,\,Y<1)=\int_0^1\left(\int_1^\infty 2e^{-x}e^{-2y}\,dx\right)dy.
	\]
	Inner integral: $y$ is held fixed, so $2e^{-2y}$ is a constant factor.
	An antiderivative of $e^{-x}$ is $-e^{-x}$, and $e^{-x}\to0$ as
	$x\to\infty$:
	\[
		2e^{-2y}\int_1^\infty e^{-x}\,dx
		=2e^{-2y}\Big[-e^{-x}\Big]_1^\infty
		=2e^{-2y}\big(0-(-e^{-1})\big)
		=2e^{-1}e^{-2y}.
	\]
	Outer integral: an antiderivative of $e^{-2y}$ is $-\frac12e^{-2y}$:
	\[
		2e^{-1}\int_0^1 e^{-2y}\,dy
		=2e^{-1}\Big[-\tfrac12e^{-2y}\Big]_0^1
		=2e^{-1}\cdot\tfrac12\big(1-e^{-2}\big)
		=e^{-1}\big(1-e^{-2}\big).
	\]
	Numerically, $e^{-1}\approx0.3679$ and $1-e^{-2}\approx0.8647$, so the
	probability is about $0.318$.

	Notice that the integral split into a product of an $x$-integral and a
	$y$-integral. This happens because the density is a product of a function
	of $x$ and a function of $y$, and the region is a rectangle.

	\step{(b)} For $x>0$, integrate out $y$ (now $x$ is held fixed):
	\[
		f_X(x)=\int_0^\infty 2e^{-x}e^{-2y}\,dy
		=2e^{-x}\Big[-\tfrac12e^{-2y}\Big]_0^\infty
		=2e^{-x}\big(0+\tfrac12\big)
		=e^{-x}.
	\]
	For $y>0$, integrate out $x$ (now $y$ is held fixed):
	\[
		f_Y(y)=\int_0^\infty 2e^{-x}e^{-2y}\,dx
		=2e^{-2y}\Big[-e^{-x}\Big]_0^\infty
		=2e^{-2y}(0+1)
		=2e^{-2y}.
	\]
	Both marginals are $0$ for negative arguments (the joint density is $0$
	there). So $X$ is exponential with rate $1$ and $Y$ is exponential with
	rate $2$. Also $f(x,y)=f_X(x)f_Y(y)$ for all $x$, $y$, so $X$ and $Y$
	are independent.

	\step{(c)} The event $X<Y$ is the part of the first quadrant above the
	line $y=x$ (shaded).
	\begin{center}
		\begin{tikzpicture}[scale=1.1,>=stealth]
			\fill[teal!25] (0,0)--(3,3)--(0,3)--cycle;
			\draw[->] (-0.2,0)--(3.3,0) node[right] {$x$};
			\draw[->] (0,-0.2)--(0,3.3) node[above] {$y$};
			\draw[navy,thick] (0,0)--(3,3) node[above right] {$y=x$};
			\draw[gold,very thick] (0,1.8)--(1.8,1.8);
			\fill[gold] (1.8,1.8) circle (1.5pt);
			\node[left] at (0,1.8) {$y$};
			\node at (0.9,2.4) {$x<y$};
			\node[below left] at (0,0) {$0$};
		\end{tikzpicture}
	\end{center}
	Fix $y>0$ (the gold segment): then $x$ runs from $0$ to $y$. So the inner
	integral is over $x$ from $0$ to $y$, and the outer one over $y$ from $0$
	to $\infty$:
	\[
		P(X<Y)=\int_0^\infty\left(\int_0^y 2e^{-x}e^{-2y}\,dx\right)dy.
	\]
	Inner integral ($y$ fixed):
	\[
		2e^{-2y}\int_0^y e^{-x}\,dx
		=2e^{-2y}\Big[-e^{-x}\Big]_0^y
		=2e^{-2y}\big(-e^{-y}-(-1)\big)
		=2e^{-2y}\big(1-e^{-y}\big)
		=2e^{-2y}-2e^{-3y}.
	\]
	Outer integral: antiderivatives are $-e^{-2y}$ and $\frac23e^{-3y}$, and
	both exponentials tend to $0$ as $y\to\infty$:
	\[
		\int_0^\infty\big(2e^{-2y}-2e^{-3y}\big)\,dy
		=\Big[-e^{-2y}+\tfrac23e^{-3y}\Big]_0^\infty
		=0-\big(-1+\tfrac23\big)
		=1-\tfrac23=\frac13.
	\]

	\begin{answer}
		(a) $P(X>1,\,Y<1)=e^{-1}(1-e^{-2})\approx0.318$ \quad
		(b) $f_X(x)=e^{-x}$ for $x>0$, $f_Y(y)=2e^{-2y}$ for $y>0$ (both $0$
		otherwise) \quad
		(c) $P(X<Y)=\frac13$
	\end{answer}
	\begin{remark}
		\emph{Interpretation.} $Y$ has the larger rate, so it tends to be
		smaller than $X$; it is no surprise that $P(X<Y)$ is below
		$\frac12$. Since $X$ and $Y$ are independent (the joint density is the
		product of the marginals),
		(a) can also be computed from the marginals:
		\[ P(X>1,\,Y<1)=P(X>1)\,P(Y<1)=e^{-1}\cdot(1-e^{-2}). \]

		\textbf{Common mistake:} in (c), writing constant bounds for both
		variables. The region $x<y$ is not a rectangle: once $y$ is fixed,
		$x$ runs only up to $y$, so the inner bound depends on the outer
		variable.
	\end{remark}
\end{solution}
